\documentclass[10pt,a4paper]{amsart}
\usepackage{amsmath,amssymb,amsthm,graphicx}
% Retained preamble for the Stacks Project chapter divisors.tex (Divisors).
% The chapter file pinned-original-source.tex carries only \input{preamble} on its first line, so no
% preamble text lives inside the chapter itself.  The Stacks Project uses one single shared preamble for
% every chapter; the snapshot stages it verbatim in sources/stacks-chow/preamble.tex.  The block below is
% copied mechanically, line for line, from that file's lines 14-18 (the Xy-pic setup) and lines 154-217
% (the native theorem declarations together with the mathematical macros).  Nothing was added, removed or
% redefined.  The wrapper class and the amsmath/amssymb/amsthm/graphicx base packages are supplied by the
% assembler, so the class/\externaldocument part of the shared preamble is not repeated here.
% For commutative diagrams we use Xy-pic
\usepackage[all]{xy}

% We use 2cell for 2-commutative diagrams.
\xyoption{2cell}
% Theorem environments.
%
\theoremstyle{plain}
\newtheorem{theorem}[subsection]{Theorem}
\newtheorem{proposition}[subsection]{Proposition}
\newtheorem{lemma}[subsection]{Lemma}

\theoremstyle{definition}
\newtheorem{definition}[subsection]{Definition}
\newtheorem{example}[subsection]{Example}
\newtheorem{exercise}[subsection]{Exercise}
\newtheorem{situation}[subsection]{Situation}

\theoremstyle{remark}
\newtheorem{remark}[subsection]{Remark}
\newtheorem{remarks}[subsection]{Remarks}

\numberwithin{equation}{subsection}

% Macros
%
\def\lim{\mathop{\mathrm{lim}}\nolimits}
\def\colim{\mathop{\mathrm{colim}}\nolimits}
\def\Spec{\mathop{\mathrm{Spec}}}
\def\Hom{\mathop{\mathrm{Hom}}\nolimits}
\def\Ext{\mathop{\mathrm{Ext}}\nolimits}
\def\SheafHom{\mathop{\mathcal{H}\!\mathit{om}}\nolimits}
\def\SheafExt{\mathop{\mathcal{E}\!\mathit{xt}}\nolimits}
\def\Sch{\mathit{Sch}}
\def\Mor{\mathop{\mathrm{Mor}}\nolimits}
\def\Ob{\mathop{\mathrm{Ob}}\nolimits}
\def\Sh{\mathop{\mathit{Sh}}\nolimits}
\def\NL{\mathop{N\!L}\nolimits}
\def\CH{\mathop{\mathrm{CH}}\nolimits}
\def\proetale{{pro\text{-}\acute{e}tale}}
\def\etale{{\acute{e}tale}}
\def\QCoh{\mathit{QCoh}}
\def\Ker{\mathop{\mathrm{Ker}}}
\def\Im{\mathop{\mathrm{Im}}}
\def\Coker{\mathop{\mathrm{Coker}}}
\def\Coim{\mathop{\mathrm{Coim}}}

% Boxtimes
%
\DeclareMathSymbol{\boxtimes}{\mathbin}{AMSa}{"02}

%
% Macros for moduli stacks/spaces
%
\def\QCohstack{\mathcal{QC}\!\mathit{oh}}
\def\Cohstack{\mathcal{C}\!\mathit{oh}}
\def\Spacesstack{\mathcal{S}\!\mathit{paces}}
\def\Quotfunctor{\mathrm{Quot}}
\def\Hilbfunctor{\mathrm{Hilb}}
\def\Curvesstack{\mathcal{C}\!\mathit{urves}}
\def\Polarizedstack{\mathcal{P}\!\mathit{olarized}}
\def\Complexesstack{\mathcal{C}\!\mathit{omplexes}}
% \Pic is the operator that assigns to X its picard group, usage \Pic(X)
% \Picardstack_{X/B} denotes the Picard stack of X over B
% \Picardfunctor_{X/B} denotes the Picard functor of X over B
\def\Pic{\mathop{\mathrm{Pic}}\nolimits}
\def\Picardstack{\mathcal{P}\!\mathit{ic}}
\def\Picardfunctor{\mathrm{Pic}}
\def\Deformationcategory{\mathcal{D}\!\mathit{ef}}
\usepackage{cleveref}
\crefname{originalreference}{original reference}{original references}
\Crefname{originalreference}{Original reference}{Original references}
\makeatletter
\newlabel{algebra-lemma-depth-ext}{{algebra-lemma-depth-ext}{source}{}{sourcepointer-1}{}}
\newlabel{algebra-lemma-depth-ext@cref}{{[originalreference][0][]algebra-lemma-depth-ext}{[1][0][]source}{}{}{}}
\makeatother
\crefname{originalreference}{original reference}{original references}
\Crefname{originalreference}{Original reference}{Original references}
\makeatletter
\newlabel{coherent-lemma-flat-base-change-cohomology}{{coherent-lemma-flat-base-change-cohomology}{source}{}{sourcepointer-2}{}}
\newlabel{coherent-lemma-flat-base-change-cohomology@cref}{{[originalreference][0][]coherent-lemma-flat-base-change-cohomology}{[1][0][]source}{}{}{}}
\makeatother
\crefname{originallemma}{original lemma}{original lemmas}
\Crefname{originallemma}{Original lemma}{Original lemmas}
\makeatletter
\newlabel{lemma-ass-weakly-ass}{{lemma-ass-weakly-ass}{source}{}{sourcepointer-3}{}}
\newlabel{lemma-ass-weakly-ass@cref}{{[originallemma][0][]lemma-ass-weakly-ass}{[1][0][]source}{}{}{}}
\makeatother
\crefname{originallemma}{original lemma}{original lemmas}
\Crefname{originallemma}{Original lemma}{Original lemmas}
\makeatletter
\newlabel{lemma-weakass-pushforward}{{lemma-weakass-pushforward}{source}{}{sourcepointer-4}{}}
\newlabel{lemma-weakass-pushforward@cref}{{[originallemma][0][]lemma-weakass-pushforward}{[1][0][]source}{}{}{}}
\makeatother
\title{Depth and nonzerodivisor statements for the pushforward of a quasi-coherent module along a quasi-compact quasi-separated morphism}
\author{The Stacks Project Authors}
\date{stacks-project:divisors.tex:local-snapshot-20261012; GNU FDL 1.2 or later}
\begin{document}
\maketitle
\noindent\textit{Editorial note.} All delivered mathematical text is the author's own, copied from the pinned Stacks Project chapter file pinned-original-source.tex; no statement or proof step was rewritten, reordered or removed.\par\medskip
\noindent\textit{Editorial note.} Following the multi-claim convention for this source, source\_url is the source-level URL https://stacks.math.columbia.edu/ and the source label lemma-depth-pushforward is carried in claim\_key and here rather than in a per-claim tag URL.\par\medskip
\noindent\textit{Editorial note.} The preamble is the chapter's shared Stacks Project preamble, retained as sources/stacks-divisors/retained-preamble.tex: a line-for-line excerpt of the shared Stacks preamble staged with this snapshot. No macro was invented or redefined, and this card adds no preamble text of its own.\par\medskip
\noindent\textit{Editorial note.} Statement span 938-961 and proof span 963-1032 come from the halfA pool card operations/pool-cards/stacks-divisors-halfA.json and were re-read in the pinned chapter file before this card was written.\par\medskip
\noindent\textit{Editorial note.} The author's own disposition of parts (1) and (2) is a one-line citation in each case and is preserved verbatim; those two parts are not counted as a proof carried out in this span.\par\medskip
\noindent\textit{Editorial note.} The proof contains a typographical slip in the author's text (a misspelling of the word determining) which was kept unchanged because the rule is to deliver the author's text as it stands.\par\medskip
\noindent\textit{Editorial note.} 4 cross-references occur inside the counted spans; they are listed in external\_references as navigation pointers only.\par\medskip
\noindent\textit{Allowed standard prerequisites (editorial).}\begin{itemize}
\item The same-chapter statements that parts (1) and (2) point at: the weak association result for pushforwards and the relation between associated and weakly associated primes. Parts (1) and (2) of the proof consist of exactly these two citations, at proof lines 964-965. They are the author's own disposition of parts (1) and (2); the counted argument therefore lives in parts (3) to (6).
\item The Ext characterisation of depth two: a local module has depth at least two exactly when both the module of homomorphisms from the residue field and its first Ext vanish, as cited from the algebra chapter. Cited at proof lines 968-970. Turns the depth statement of part (3) into the vanishing of one homomorphism group and one Ext group.
\item Flat base change for the cohomology of quasi-coherent modules, as cited from the chapter on coherent cohomology. Cited at proof lines 976-977 and again at 1000-1002. Lets the proof replace the base by the spectrum of the local ring at the point and the source by the corresponding fibre product.
\item Adjunction between pushforward and pullback for quasi-coherent modules. Used at proof lines 987-989 in the identification of homomorphism groups over S with homomorphism groups over X. Moves the homomorphism group along f to the source, where the hypothesis that the point is not in the image makes the pulled-back maximal ideal the whole structure sheaf.
\end{itemize}\medskip
\ifdefined\hypertarget\hypertarget{sourcepointer-1}{}\fi
\noindent\textit{Source reference.} Original reference algebra-lemma-depth-ext: sources/stacks-divisors/pinned-original-source.tex, source lines 971--971. Same-chapter reference used in the retained author text at source line(s) 971. It is supplied only as a navigation pointer so that the retained text compiles; the counted item delivers the author's own statement and the author's complete proof, and no part of the referenced item is copied in or substituted for the counted argument.\par\medskip
\ifdefined\hypertarget\hypertarget{sourcepointer-2}{}\fi
\noindent\textit{Source reference.} Original reference coherent-lemma-flat-base-change-cohomology: sources/stacks-divisors/pinned-original-source.tex, source lines 737--737. Same-chapter reference used in the retained author text at source line(s) 737, 981, 1000, 7609. It is supplied only as a navigation pointer so that the retained text compiles; the counted item delivers the author's own statement and the author's complete proof, and no part of the referenced item is copied in or substituted for the counted argument.\par\medskip
\ifdefined\hypertarget\hypertarget{sourcepointer-3}{}\fi
\noindent\textit{Source reference.} Original lemma lemma-ass-weakly-ass: sources/stacks-divisors/pinned-original-source.tex, source lines 701--701. Same-chapter reference used in the retained author text at source line(s) 875, 965, 1222, 2625, 6026, 6345. It is supplied only as a navigation pointer so that the retained text compiles; the counted item delivers the author's own statement and the author's complete proof, and no part of the referenced item is copied in or substituted for the counted argument.\par\medskip
\ifdefined\hypertarget\hypertarget{sourcepointer-4}{}\fi
\noindent\textit{Source reference.} Original lemma lemma-weakass-pushforward: sources/stacks-divisors/pinned-original-source.tex, source lines 720--720. Same-chapter reference used in the retained author text at source line(s) 793, 964. It is supplied only as a navigation pointer so that the retained text compiles; the counted item delivers the author's own statement and the author's complete proof, and no part of the referenced item is copied in or substituted for the counted argument.\par\medskip
\begin{lemma}
\label{lemma-depth-pushforward}
Let $f : X \to S$ be a quasi-compact and quasi-separated morphism of schemes.
Let $\mathcal{F}$ be a quasi-coherent $\mathcal{O}_X$-module.
Let $s \in S$.
\begin{enumerate}
\item If $s \not \in f(X)$, then $s$ is not weakly associated
to $f_*\mathcal{F}$.
\item If $s \not \in f(X)$ and $\mathcal{O}_{S, s}$ is Noetherian,
then $s$ is not associated to $f_*\mathcal{F}$.
\item If $s \not \in f(X)$, $(f_*\mathcal{F})_s$ is a finite
$\mathcal{O}_{S, s}$-module, and $\mathcal{O}_{S, s}$
is Noetherian, then $\text{depth}((f_*\mathcal{F})_s) \geq 2$.
\item If $\mathcal{F}$ is flat over $S$ and $a \in \mathfrak m_s$
is a nonzerodivisor, then $a$ is a nonzerodivisor on $(f_*\mathcal{F})_s$.
\item If $\mathcal{F}$ is flat over $S$ and $a, b \in \mathfrak m_s$
is a regular sequence, then $a$ is a nonzerodivisor on $(f_*\mathcal{F})_s$
and $b$ is a nonzerodivisor on $(f_*\mathcal{F})_s/a(f_*\mathcal{F})_s$.
\item If $\mathcal{F}$ is flat over $S$ and $(f_*\mathcal{F})_s$
is a finite $\mathcal{O}_{S, s}$-module, then
$\text{depth}((f_*\mathcal{F})_s) \geq
\min(2, \text{depth}(\mathcal{O}_{S, s}))$.
\end{enumerate}
\end{lemma}

\begin{proof}
Part (1) is Lemma \ref{lemma-weakass-pushforward}.
Part (2) follows from (1) and Lemma \ref{lemma-ass-weakly-ass}.

\medskip\noindent
Proof of part (3). To show the depth is $\geq 2$ it suffices to show that
$\Hom_{\mathcal{O}_{S, s}}(\kappa(s), (f_*\mathcal{F})_s) = 0$ and
$\Ext^1_{\mathcal{O}_{S, s}}(\kappa(s), (f_*\mathcal{F})_s) = 0$, see
Algebra, Lemma \ref{algebra-lemma-depth-ext}.
Using the exact sequence
$0 \to \mathfrak m_s \to \mathcal{O}_{S, s} \to \kappa(s) \to 0$
it suffices to prove that the map
$$
\Hom_{\mathcal{O}_{S, s}}(\mathcal{O}_{S, s}, (f_*\mathcal{F})_s)
\to
\Hom_{\mathcal{O}_{S, s}}(\mathfrak m_s, (f_*\mathcal{F})_s)
$$
is an isomorphism. By flat base change (Cohomology of Schemes, Lemma
\ref{coherent-lemma-flat-base-change-cohomology})
we may replace $S$ by
$\Spec(\mathcal{O}_{S, s})$ and $X$ by $\Spec(\mathcal{O}_{S, s}) \times_S X$.
Denote $\mathfrak m \subset \mathcal{O}_S$ the ideal sheaf of $s$.
Then we see that
$$
\Hom_{\mathcal{O}_{S, s}}(\mathfrak m_s, (f_*\mathcal{F})_s) =
\Hom_{\mathcal{O}_S}(\mathfrak m, f_*\mathcal{F}) =
\Hom_{\mathcal{O}_X}(f^*\mathfrak m, \mathcal{F})
$$
the first equality because $S$ is local with closed point $s$
and the second equality
by adjunction for $f^*, f_*$ on quasi-coherent modules. However, since
$s \not \in f(X)$ we see that $f^*\mathfrak m = \mathcal{O}_X$.
Working backwards through the arguments we get the desired equality.

\medskip\noindent
For the proof of (4), (5), and (6) we use flat base change
(Cohomology of Schemes, Lemma
\ref{coherent-lemma-flat-base-change-cohomology})
to reduce to the case where $S$ is the spectrum of
$\mathcal{O}_{S, s}$.
Then a nonzerodivisor $a \in \mathcal{O}_{S, s}$
deterimines a short exact sequence
$$
0 \to \mathcal{O}_S \xrightarrow{a} \mathcal{O}_S \to
\mathcal{O}_S/a \mathcal{O}_S \to 0
$$
Since $\mathcal{F}$ is flat over $S$, we obtain an exact sequence
$$
0 \to \mathcal{F} \xrightarrow{a} \mathcal{F} \to
\mathcal{F}/a\mathcal{F} \to 0
$$
Pushing forward we obtain an exact sequence
$$
0 \to f_*\mathcal{F} \xrightarrow{a} f_*\mathcal{F} \to
f_*(\mathcal{F}/a\mathcal{F})
$$
This proves (4) and it shows that
$f_*\mathcal{F}/ af_*\mathcal{F} \subset f_*(\mathcal{F}/a\mathcal{F})$.
If $b$ is a nonzerodivisor on
$\mathcal{O}_{S, s}/a\mathcal{O}_{S, s}$, then the exact same argument shows
$b : \mathcal{F}/a\mathcal{F} \to \mathcal{F}/a\mathcal{F}$
is injective. Pushing forward we conclude
$$
b : f_*(\mathcal{F}/a\mathcal{F}) \to f_*(\mathcal{F}/a\mathcal{F})
$$
is injective and hence also
$b : f_*\mathcal{F}/ af_*\mathcal{F} \to f_*\mathcal{F}/ af_*\mathcal{F}$
is injective. This proves (5). Part (6) follows from
(4) and (5) and the definitions.
\end{proof}

\end{document}
